\ifdefined\gkver\else\def\gkver{student}\fi
\documentclass[\gkver]{gaokaozhenti}
\begin{document}

\chapter{1987年全国}

\begin{enumerate}
\item
%% number: 1
%% paperName: 1987年全国高考第 1 题 3 分
%% typeId: 3
%% type: 填空题
%% chapter: 原子和原子核
%% point: 玻尔模型
%% method: 无
%% score: 3
%% degree: 650
%% duplicateId: 0
%% body:
在玻尔的氢原子模型中，电子的第一条（即离核最近的那条）可能轨道的半径为 $r_{1}$，则由此向外数的第三条可能轨道的半径 $r_{3}=$ \tkanswer[2cm]{$9r_{1}$}，电子在这第三条可能轨道上运动时的动能 $E_{k}=$ \tkanswer[2.6cm]{$\frac{ke^{2}}{18r_{1}}$}。已知基本电荷为 $e$，静电力恒量为 $k$。

%% answer: 9r_{1}，ke^{2}/18r_{1}
\jdanswer{
$9r_{1}$，$\frac{ke^{2}}{18r_{1}}$
}

%% memo:
\memoanswer{
由 $r_{n}=n^{2}r_{1}$ 知 $r_{3}=9r_{1}$，由库仑力提供向心力得 $\frac{ke^{2}}{r_{3}^{2}}=m\frac{v^{2}}{r_{3}}$，故 $E_{k3}=\frac{1}{2}mv^{2}=\frac{ke^{2}}{18r_{1}}$。

评分标准：全题 24 分，每小题 3 分。
}

\item
%% number: 2
%% paperName: 1987年全国高考第 2 题 3 分
%% typeId: 3
%% type: 填空题
%% chapter: 原子和原子核
%% point: 人工核反应
%% method: 无
%% score: 3
%% degree: 650
%% duplicateId: 0
%% body:
一个锂核（\ce{^{7}_{3}Li}）受到一个质子的轰击，变成两个 $\alpha$ 粒子，这一过程的核反应方程是 \tkanswer[5cm]{\ce{^{7}_{3}Li{}+{}^{1}_{1}H -> 2^{4}_{2}He}}。已知一个氢原子的质量是 $1.6736\times10^{-27}\Ukg$，一个锂原子的质量是 $11.6505\times10^{-27}\Ukg$，一个氦原子的质量是 $6.6466\times10^{-27}\Ukg$。上述核反应所释放的能量等于 \tkanswer[3cm]{$2.78\times10^{-12}\UJ$}。（最后结果取三位有效数字。）

%% answer: ^{7}_{3}Li + ^{1}_{1}H -> 2^{4}_{2}He，2.78×10^{-12}
\jdanswer{
$\ce{^{7}_{3}Li + ^{1}_{1}H -> 2^{4}_{2}He}$，$2.78\times10^{-12}\UJ$
}

%% memo:
\memoanswer{
质量亏损为 $\Delta m=11.6505\times10^{-27}+1.6736\times10^{-27}-2\times6.6466\times10^{-27}=0.0311\times10^{-27}\Ukg$，

由质能方程得 $\Delta E=\Delta mc^{2}=0.0311\times10^{-27}\times(3\times10^{8})^{2}=2.78\times10^{-12}\UJ$。
}

\item
%% number: 3
%% paperName: 1987年全国高考第 3 题 3 分
%% typeId: 3
%% type: 填空题
%% chapter: 光学
%% point: 光的干涉
%% method: 无
%% score: 3
%% degree: 750
%% duplicateId: 0
%% body:
在太阳光照射下，水面油膜上会出现彩色的花纹，这是两列相干光波发生干涉的结果。这两列相干光波是太阳光分别经 \tkanswer[4cm]{油膜的前表面和后表面反射} 而形成的。用平行的单色光垂直照射不透明的小圆板，在圆板后面的屏上发现圆板阴影中心处有一个亮斑，这是光的 \tkanswer{衍射} 现象。

%% answer: 油膜的前表面和后表面反射，衍射
\jdanswer{
油膜的前表面和后表面反射，衍射
}

%% memo:
\memoanswer{
用平行的单色光垂直照射不透明的小圆板，在圆板后面的屏上发现圆板阴影中心处有一个亮斑，就是著名的“泊松亮斑”，是衍射实验的结果。

评分标准：答“油膜的前表面和后表面反射”给 2 分，答“衍射”给 1 分。
}

\item
%% number: 4
%% paperName: 1987年全国高考第 4 题 3 分
%% typeId: 3
%% type: 填空题
%% chapter: 力物体的平衡
%% point: 摩擦力
%% method: 无
%% score: 3
%% degree: 650
%% duplicateId: 0
%% body:
一根质量为 $m$，长度为 $l$ 的均匀的长方木料放在水平桌面上，木料与桌面间的摩擦系数为 $\mu$。现用水平力 $F$ 推木料，当木料经过图中所示的位置时，桌面对它的摩擦力等于 \tkanswer{$\mu mg$}。
\onepicture[align=right, w=6cm]{figs/04.png}

%% answer: \mu mg
\jdanswer{
$\mu mg$
}

%% memo:
\memoanswer{
因木料的重心在水平桌面上，所以压力 $N=mg$，所以摩擦力 $f=\mu mg$。
}

\item
%% number: 5
%% paperName: 1987年全国高考第 5 题 3 分
%% typeId: 7
%% type: 作图题
%% chapter: 机械波
%% point: 波的图象
%% method: 无
%% score: 3
%% degree: 750
%% duplicateId: 0
%% body:
绳中有一列正弦横波，沿 $x$ 轴传播。右上图中 a、b 是绳上两点，它们在 $x$ 轴方向上的距离小于一个波长，当 a 点振动到最高点时，b 点恰经过平衡位置向上运动。试在图上 a、b 之间画出两个波形分别表示：①沿 $x$ 轴正方向传播的波；②沿 $x$ 轴负方向传播的波。在所画波形上要注明符号①和②。
\onepicture[align=right, w=6cm]{\includesvg{figs/05}}

%% answer: 如图所示
\jdanswer{
如图所示
}

%% memo:
\memoanswer{
若波沿 $x$ 轴正方向传播，则因 b 点恰经过平衡位置向上运动，所以处在从左往右看的“下坡”位置，如图中①；若波沿 $x$ 轴负方向传播，则因 b 点恰经过平衡位置向上运动，所以处在从右往左看的“下坡”位置，如图中②。画图时先画草图再往试卷上画。

\onepicture[w=6cm]{\includesvg{figs/05_an}}

评分标准：全题 3 分，画对一个波形的，只给 1 分。
}

\item
%% number: 6
%% paperName: 1987年全国高考第 6 题 3 分
%% typeId: 3
%% type: 填空题
%% chapter: 电场
%% point: 带电粒子的圆周运动
%% method: 等效替代
%% score: 3
%% degree: 450
%% duplicateId: 0
%% body:
电子绕核运动可等效为一环形电流。设氢原子中的电子以速率 $v$ 在半径为 $r$ 的轨道上运动，用 $e$ 表示电子的电量，则其等效电流的电流强度等于 \tkanswer[2.6cm]{$\frac{ev}{2\pi r}$}。

%% answer: \frac{v}{2\pi r}e
\jdanswer{
$\frac{ev}{2\pi r}$
}

%% memo:
\memoanswer{
由 $I=\frac{q}{t}$ 得 $I=\frac{e}{T}$，因 $T=\frac{2\pi r}{v}$，故 $I=\frac{ev}{2\pi r}$。
}

\item
%% number: 7
%% paperName: 1987年全国高考第 7 题 3 分
%% typeId: 3
%% type: 填空题
%% chapter: 交流电
%% point: 变压器
%% method: 无
%% score: 3
%% degree: 550
%% duplicateId: 0
%% body:
如图是一台三相变压器的示意图，原线圈和副线圈的匝数比为 $n_{1}/n_{2}=30$，今测得副线圈一边 a、b、c 三端的任意两端之间的电压为 $380\UV$，则在原线圈一边 A、B、C 三端的任意两端之间的电压等于 \tkanswer[2cm]{6600} $\UV$。
\onepicture[align=right, w=7cm]{figs/07.png}

%% answer: 6600
\jdanswer{
$6600\UV$
}

%% memo:
\memoanswer{
因为副线圈一边是星形接法，$U_{\text{线}}=\sqrt{3}U_{\text{相}}$，已知 $U_{\text{线}}=380\UV$，所以 $U_{\text{相}}=220\UV$，由变压器原理得 $\frac{U_{\text{原}}}{U_{\text{副}}}=\frac{n_{1}}{n_{2}}$，所以 $U_{\text{原}}=220\times30\UV=6600\UV$，因为原线圈一边是三角形接法，所以 $U_{\text{线}}=U_{\text{相}}=6600\UV$。
}

\item
%% number: 8
%% paperName: 1987年全国高考第 8 题 3 分
%% typeId: 7
%% type: 作图题
%% chapter: 气体性质
%% point: 气体图象
%% method: 无
%% score: 3
%% degree: 550
%% duplicateId: 0
%% body:
图 1 中的实线表示 1 摩尔的理想气体发生状态变化时的 $p-V$ 图线，变化过程是由状态 A 出发，经过 B、C、D 诸状态，最后又回到状态 A，试将这全部过程准确地画在图 2 所示的 $p-T$ 图中。
\onepicture[align=right, w=10cm]{figs/08.png}

%% answer: 如图所示
\jdanswer{
如图所示
}

%% memo:
\memoanswer{
由 $p-V$ 图像知，$A\to B$ 为等压变化，$B\to C$ 为等容变化，$C\to D$ 为等温变化，$D\to A$ 为等容变化。

A 点为一个标准大气压，体积为 $22.4\UL$，则温度为 $0\UdC$，即 $273\UK$。

$A\to B$ 为等压变化，有 $\frac{V_{A}}{T_{A}}=\frac{V_{B}}{T_{B}}$，则 $T_{B}=2\times273\UK$，则 B 的状态为 $(2\times273, 1)$；

$B\to C$ 为等容变化，有 $\frac{p_{B}}{T_{B}}=\frac{p_{C}}{T_{C}}$，则 $T_{C}=4\times273\UK$，则 C 的状态为 $(4\times273, 2)$；

$C\to D$ 为等温变化，$T_{D}=4\times273\UK$，则 D 的状态为 $(4\times273, 4)$；

在 $p-T$ 图像中将全部过程准确地画出来，如图所示。

\onepicture[w=6.5cm]{figs/08_an.png}
}

\item
%% number: 9
%% paperName: 1987年全国高考第 9 题 4 分
%% typeId: 2
%% type: 多选题
%% chapter: 磁场
%% point: 带电粒子在磁场中的运动
%% method: 无
%% score: 4
%% degree: 450
%% duplicateId: 0
%% body:
如图所示，连接平行金属板 $P_{1}$ 和 $P_{2}$（板面垂直于纸面）的导线的一部分 CD 和另一连接电池的回路的一部分 GH 平行，CD 和 GH 均在纸平面内，金属板置于磁场中，磁场方向垂直于纸面向里，当一束等离子体射入两金属板之间时，CD 段导线将受到力的作用。\xzanswer{AD}
\onepicture[w=6cm]{figs/09.png}
\fourchoices[answer=AD]
{等离子体从右方射入时，CD 受力的方向背离 GH}
{等离子体从右方射入时，CD 受力的方向指向 GH}
{等离子体从左方射入时，CD 受力的方向背离 GH}
{等离子体从左方射入时，CD 受力的方向指向 GH}

%% answer: AD
%% memo:
\memoanswer{
等离子体从右方射入时，正电荷受洛伦兹力向上偏转，负电荷受洛伦兹力向下偏转，所以 $P_{1}$ 带正电，$P_{2}$ 带负电，所以 CD 中电流从 C 向 D，因 GH 中电流从 G 向 H，由右手螺旋定则知，GH 中电流的磁场在 CD 处向里，由左手定则知，CD 受力的方向背离 GH，故 A 正确 B 错误；同样的方法可以判断，等离子体从左方射入时，CD 受力的方向指向 GH，故 D 正确 C 错误。
}

\item
%% number: 10
%% paperName: 1987年全国高考第 10 题 4 分
%% typeId: 1
%% type: 单选题
%% chapter: 力物体的平衡
%% point: 力矩平衡
%% method: 无
%% score: 4
%% degree: 450
%% duplicateId: 0
%% body:
某同学用一不等臂天平称量物体 A 的质量 $M$。他先把物体 A 放在天平的右方托盘上，使天平平衡时，左方托盘上所放砝码的质量为 $m_{1}$；他再把物体 A 放在天平的左方托盘上，使天平平衡时，右方托盘上所放砝码的质量为 $m_{2}$。被称物体的质量 $M$\xzanswer{A}
\fourchoices[answer=A]
{等于 $\sqrt{m_{1}m_{2}}$}
{等于 $\frac{m_{1}+m_{2}}{2}$}
{等于 $\frac{m_{1}m_{2}}{m_{1}+m_{2}}$}
{无法确定，因为所用天平是不等臂的}

%% answer: A
%% memo:
\memoanswer{
由天平的平衡条件得 $MgL_{\text{右}}=m_{1}gL_{\text{左}}$，$m_{2}gL_{\text{右}}=MgL_{\text{左}}$，联立解得 $M=\sqrt{m_{1}m_{2}}$，故 A 正确。
}

\item
%% number: 11
%% paperName: 1987年全国高考第 11 题 4 分
%% typeId: 1
%% type: 单选题
%% chapter: 电磁感应
%% point: 楞次定律
%% method: 无
%% score: 4
%% degree: 550
%% duplicateId: 0
%% body:
如图，有一固定的超导体圆环，在其右侧放着一条形磁铁，此时圆环中没有电流。当把磁铁向右方移走时，由于电磁感应，在超导体圆环中产生了一定的电流\xzanswer{D}
\onepicture[w=5cm]{figs/11.png}
\fourchoices[answer=D]
{这电流的方向如图中箭头所示，磁铁移走后，这电流很快消失}
{这电流的方向如图中箭头所示，磁铁移走后，这电流继续维持}
{这电流的方向与图中箭头方向相反，磁铁移走后，这电流很快消失}
{这电流的方向与图中箭头方向相反，磁铁移走后，这电流继续维持}

%% answer: D
%% memo:
\memoanswer{
由楞次定律知，感应电流的磁场方向向左，所以感应电流的方向与图中箭头方向相反，因为超导体的电阻为 0，所以磁铁移走后，这电流继续维持，故 D 正确。
}

\item
%% number: 12
%% paperName: 1987年全国高考第 12 题 4 分
%% typeId: 2
%% type: 多选题
%% chapter: 曲线运动
%% point: 人造地球卫星
%% method: 无
%% score: 4
%% degree: 350
%% duplicateId: 0
%% body:
用 $m$ 表示地球通讯卫星（同步卫星）的质量，$h$ 表示它离地面的高度，$R_{0}$ 表示地球的半径，$g_{0}$ 表示地球表面处的重力加速度，$\omega_{0}$ 表示地球自转的角速度，则通讯卫星所受的地球对它的万有引力的大小\xzanswer{BC}
\fourchoices[answer=BC]
{等于 0}
{等于 $m\frac{R_{0}^{2}g_{0}}{(R_{0}+h)^{2}}$}
{等于 $m\sqrt[3]{R_{0}^{2}g_{0}\omega_{0}^{4}}$}
{以上结果都不正确}

%% answer: BC
%% memo:
\memoanswer{
地球表面处 $G\frac{Mm}{R_{0}^{2}}=mg_{0}$，得 $GM=g_{0}R_{0}^{2}$，通讯卫星所受的地球对它的万有引力的大小 $F=G\frac{Mm}{(R_{0}+h)^{2}}=\frac{mg_{0}R_{0}^{2}}{(R_{0}+h)^{2}}$，故 B 正确；由万有引力提供向心力得 $G\frac{Mm}{(R_{0}+h)^{2}}=m\omega_{0}^{2}(R_{0}+h)$，得 $R_{0}+h=\sqrt[3]{\frac{GM}{\omega_{0}^{2}}}$，所以通讯卫星所受的万有引力大小 $F=m\omega_{0}^{2}(R_{0}+h)=m\sqrt[3]{R_{0}^{2}g_{0}\omega_{0}^{4}}$，故 C 正确。
}

\item
%% number: 13
%% paperName: 1987年全国高考第 13 题 4 分
%% typeId: 1
%% type: 单选题
%% chapter: 电场
%% point: 电场叠加
%% method: 无
%% score: 4
%% degree: 550
%% duplicateId: 0
%% body:
在 $x$ 轴上有两个点电荷，一个带正电 $Q_{1}$，一个带负电 $-Q_{2}$，且 $Q_{1}=2Q_{2}$。用 $E_{1}$ 和 $E_{2}$ 分别表示两个电荷所产生的场强的大小，则在 $x$ 轴上\xzanswer{B}
\fourchoices[answer=B]
{$E_{1}=E_{2}$ 之点只有一处；该处合场强为 0}
{$E_{1}=E_{2}$ 之点共有两处；一处合场强为 0，另一处合场强为 $2E_{2}$}
{$E_{1}=E_{2}$ 之点共有三处；其中两处合场强为 0，另一处合场强为 $2E_{2}$}
{$E_{1}=E_{2}$ 之点共有三处；其中一处合场强为 0，另两处合场强为 $2E_{2}$}

%% answer: B
%% memo:
\memoanswer{
\onepicture[w=8cm]{figs/13_an.png}

如图所示，设两个点电荷之间的距离为 $l$，某点距 $Q_{1}$ 为 $x$，则距 $Q_{2}$ 为 $l-x$（当 $x<l$）或 $x-l$（当 $x>l$），则 $E_{1}=\frac{2kQ}{x^{2}}$，$E_{2}=\frac{kQ}{(l-x)^{2}}$ 或 $E_{2}=\frac{kQ}{(x-l)^{2}}$，令 $E_{1}=E_{2}$，解得 $x=\frac{\sqrt{2}}{\sqrt{2}+1}l$ 或 $x=\frac{\sqrt{2}}{\sqrt{2}-1}l$，所以 $E_{1}=E_{2}$ 的点共有两处；一处合场强为 0，另一处合场强为 $2E_{2}$，故 B 正确。
}

\item
%% number: 14
%% paperName: 1987年全国高考第 14 题 4 分
%% typeId: 1
%% type: 单选题
%% chapter: 电磁感应
%% point: 自感与互感，涡流
%% method: 无
%% score: 4
%% degree: 650
%% duplicateId: 0
%% body:
如图所示电路，多匝线圈的电阻和电池的内电阻可以忽略，两个电阻器的阻值都是 $R$。电键 S 原来打开着，电流 $I_{0}=\frac{E}{2R}$。今合下电键将一个电阻器短路，于是线圈中有自感电动势产生，这自感电动势\xzanswer{D}
\onepicture[w=5.5cm]{figs/14.png}
\fourchoices[answer=D]
{有阻碍电流的作用，最后电流由 $I_{0}$ 减小为零}
{有阻碍电流的作用，最后电流总小于 $I_{0}$}
{有阻碍电流增大的作用，因而电流保持为 $I_{0}$ 不变}
{有阻碍电流增大的作用，但电流最后还是要增大到 $2I_{0}$}

%% answer: D
%% memo:
\memoanswer{
线圈产生的自感电流阻碍原电流的变化，闭合开关，电路的电流增大，故线圈有阻碍电流增大的作用，阻碍不是阻止，线圈最后通过的稳定电流为 $I=\frac{E}{R}=2I_{0}$，故 D 正确。
}

\item
%% number: 15
%% paperName: 1987年全国高考第 15 题 4 分
%% typeId: 1
%% type: 单选题
%% chapter: 运动和力
%% point: 连接体
%% method: 无
%% score: 4
%% degree: 550
%% duplicateId: 0
%% body:
如图所示，一根轻质弹簧上端固定，下端挂一质量为 $m_{0}$ 的平盘，盘中有一物体，质量为 $m$。当盘静止时，弹簧的长度比其自然长度伸长了 $l$。今向下拉盘使弹簧再伸长 $\Delta l$ 后停止，然后松手放开。设弹簧总处在弹性限度以内，则刚松开手时盘对物体的支持力等于\xzanswer{A}
\onepicture[w=2.2cm]{figs/15.png}
\fourchoices[answer=A]
{$\left(1+\frac{\Delta l}{l}\right)mg$}
{$\left(1+\frac{\Delta l}{l}\right)(m+m_{0})g$}
{$\frac{\Delta l}{l}mg$}
{$\frac{\Delta l}{l}(m+m_{0})g$}

%% answer: A
%% memo:
\memoanswer{
把盘和物体看做整体，对系统整体分析，由平衡条件得 $kl=(m+m_{0})g$，由牛顿第二定律得 $k\Delta l=(m+m_{0})a$；隔离物体，由牛顿第二定律有 $F_{N}-mg=ma$，联立解得 $F_{N}=mg\left(1+\frac{\Delta l}{l}\right)$，故 A 正确。
}

\item
%% number: 16
%% paperName: 1987年全国高考第 16 题 4 分
%% typeId: 2
%% type: 多选题
%% chapter: 光学
%% point: 透镜
%% method: 无
%% score: 4
%% degree: 650
%% duplicateId: 0
%% body:
下图中的 $L_{1}$、$L_{2}$、$L_{3}$ 和 $L_{4}$ 分别表示放在空气中的薄透镜，$OO'$ 表示主轴，透镜的焦点没有画出，也不知是凸透镜还是凹透镜。对每个透镜，图中给定了两条入射光线。关于出射光线的画法，在这四个光路图中，哪个或哪几个光路图是近似正确的？\xzanswer{BD}
\fourchoices[answer=BD, ispicture=true, h=2.3cm]
{\includesvg{figs/16a}}
{\includesvg{figs/16b}}
{\includesvg{figs/16c}}
{\includesvg{figs/16d}}

%% answer: BD
%% memo:
\memoanswer{
图 A，一条光线折向主轴，另一条远离主光轴，互相矛盾，故 A 错误；图 B、C，虽然两条光线都远离主光轴，B 中两折射光线的延长线与副光轴的交点连线垂直于主轴，而 C 中两折射光线的延长线与副光轴的交点连线不垂直于主光轴，故 B 正确 C 错误；图 D，两条光线都折向主光轴，$L_{4}$ 是凸透镜时是可能的，故 D 正确。
}

\item
%% number: 17
%% paperName: 1987年全国高考第 17 题 4 分
%% typeId: 1
%% type: 单选题
%% chapter: 气体性质
%% point: 阿伏伽德罗常数
%% method: 无
%% score: 4
%% degree: 550
%% duplicateId: 0
%% body:
如图所示，食盐（NaCl）的晶体是由钠离子（图中 $\bigcirc$）和氯离子（图中 $\bullet$）组成的。这两种离子在空间中三个互相垂直的方向上，都是等距离地交错排列的。已知食盐的摩尔质量是 $58.5\Ugmol$，食盐的密度是 $2.2\Ugcmc$。阿伏伽德罗常数为 $6.0\times10^{23}\Umoln$。在食盐晶体中两个距离最近的钠离子中心间的距离的数值最接近于（就下面四个数值相比）\xzanswer{C}
\onepicture[w=5cm]{figs/17.png}
\fourchoices[answer=C]
{$3.0\times10^{-8}\Ucm$}
{$3.5\times10^{-8}\Ucm$}
{$4.0\times10^{-8}\Ucm$}
{$5.0\times10^{-8}\Ucm$}

%% answer: C
%% memo:
\memoanswer{
每个离子占据的空间体积为

$V=\frac{M}{\rho N_{\mathrm{A}}}=\frac{58.5\times10^{-3}}{2.2\times10^{3}\times6.0\times10^{23}}\Umc=4.43\times10^{-29}\Umc$，

相邻粒子之间的距离为

$d=\sqrt[3]{V}=\sqrt[3]{44.3\times10^{-30}}\Um=3.5\times10^{-10}\Um$，

两个距离最近的钠离子中心间的距离

$s=\sqrt{2}d=4.2\times10^{-10}\Um\approx4.0\times10^{-8}\Ucm$，故 C 正确。
}

\item
%% number: 18
%% paperName: 1987年全国高考第 18 题 4 分
%% typeId: 2
%% type: 多选题
%% chapter: 曲线运动
%% point: 圆周运动的应用
%% method: 无
%% score: 4
%% degree: 250
%% duplicateId: 0
%% body:
图中 M、N 是两个共轴圆筒的横截面，外筒半径为 $R$，内筒半径比 $R$ 小很多，可以忽略不计，筒的两端是封闭的，两筒之间抽成真空。两筒以相同的角速度 $\omega$ 绕其中心轴线（图中垂直于纸面）做匀速转动。设从 M 筒内部可以通过窄缝 S（与 M 筒的轴线平行）不断地向外射出两种不同速率 $v_{1}$ 和 $v_{2}$ 的微粒，从 S 处射出时的初速度的方向都是沿筒的半径方向，微粒到达 N 筒后就附着在 N 筒上。如果 $R$、$v_{1}$ 和 $v_{2}$ 都不变，而 $\omega$ 取某一合适的值，则\xzanswer{ABC}
\onepicture[w=5cm]{figs/18.png}
\fourchoices[answer=ABC]
{有可能使微粒落在 N 筒上的位置都在 a 处一条与 S 缝平行的窄条上}
{有可能使微粒落在 N 筒上的位置都在某一处如 b 处一条与 S 缝平行的窄条上}
{有可能使微粒落在 N 筒上的位置分别在某两处如 b 处和 c 处与 S 缝平行的窄条上}
{只要时间足够长，N 筒上将到处都落有微粒}

%% answer: ABC
%% memo:
\memoanswer{
由于 N 筒与 M 筒一起以同一角速度绕同一中心轴转动，所以在任意一段时间内，N 筒上与 S 缝在同一径向连线上的 a 点（及线）必定与 S 缝一起转过同样大的角度；因内筒 M 的半径可以忽略不计，所以，粒子射出 S 后的方向都沿半径方向；如果射出 S 缝的粒子的速度为 $v$，则从射出到到达 N 筒的时间为 $\frac{R}{v}$。

如果半径与角速度的关系为：两种粒子到达 a 点的时间差为周期的整数倍，则 A 正确；

如果半径与角速度的关系为：两种粒子到达 b 点的时间差为周期的整数倍，则 B 正确；

如果半径与角速度的关系为：两种粒子到达 a 点的时间差为周期的整数倍并相差从 b 到 c 的时间，则 C 正确。
}

\item
%% number: 19
%% paperName: 1987年全国高考第 19 题 4 分
%% typeId: 4
%% type: 实验题
%% chapter: 运动和力
%% point: 验证牛顿第二定律
%% method: 无
%% score: 4
%% degree: 850
%% duplicateId: 0
%% body:
用图甲所示的装置研究质量一定时加速度与作用力的关系。研究的对象是放在长木板上的小车，小车的质量为 $M$，长木板是水平放置的。小车前端拴着细轻绳，跨过定滑轮，下面吊着砂桶。实验中认为细绳对小车的作用力 $F$ 等于砂和桶的总重量 $mg$。用改变砂的质量的办法来改变小车的作用力 $F$，用打点计时器测出小车的加速度 $a$，得出若干组 $F$ 和 $a$ 的数据。然后根据测得的数据作出 $a-F$ 图线。

一学生作出如图乙所示的图线，发现横轴上的截距 OA 较大，明显地超出了偶然误差的范围，这是由于在实验中没有进行下面的步骤，即 \tkanswer{平衡摩擦力}。
\onepicture[align=right, w=10cm]{figs/19.png}

%% answer: 平衡摩擦力
\jdanswer{
平衡摩擦力
}

%% memo:
\memoanswer{
从图乙看出，当 $F=A$ 时加速度才开始从 0 增大，说明摩擦力不能忽略，所以要平衡摩擦力才能使图象过原点，为正比例图象。

评分标准：平衡摩擦力（4 分）。
}

\item
%% number: 20
%% paperName: 1987年全国高考第 20 题 4 分
%% typeId: 4
%% type: 实验题
%% chapter: 电路
%% point: 多用表的使用
%% method: 无
%% score: 4
%% degree: 750
%% duplicateId: 0
%% body:
一学生使用万用表测电阻，他在实验中有违反使用规则之处，他的主要实验步骤如下：
\begin{steps}
	\item
	把选择开关扳到“$\times1$K”的欧姆挡上；
	\item
	把表笔插入测试笔插孔中，先把两根表笔相接触，旋转调零旋钮，使指针指在电阻刻度的零位上；
	\item
	把两根表笔分别与某一待测电阻的两端相接，发现这时指针偏转较小；
	\item
	换用“$\times100$”的欧姆挡，发现这时指针偏转适中．随即记下欧姆数值；
	\item
	把表笔从测试笔插孔中拔出后，就把万用表放回桌上原处，实验完毕。
\end{steps}

这个学生在测量时已注意到：待测电阻与其他元件和电源断开，不用手碰表笔的金属杆。这个学生在实验中违反了哪一或哪些重要的使用规则？

答：\tkanswer[9cm]{换用欧姆挡的量程时，要重新调整零旋钮；万用表使用后不能把选择开关置于欧姆挡}。

%% answer: 换用欧姆挡的量程时，要重新调整零旋钮；万用表使用后不能把选择开关置于欧姆挡
\jdanswer{
\begin{enumerate}
	\item
	换用欧姆挡的量程时，要重新调整零旋钮。

	\item
	万用表使用后不能把选择开关置于欧姆挡。
\end{enumerate}
}

%% memo:
\memoanswer{
要注意：换用欧姆挡的量程时，要重新调整零旋钮；实验结束，需将挡位置于 OFF 挡位。

评分标准：换用欧姆挡的量程时，要重新调整零旋钮（2 分）；万用表使用后不能把选择开关置于欧姆挡（2 分）。
}

\item
%% number: 21
%% paperName: 1987年全国高考第 21 题 4 分
%% typeId: 4
%% type: 实验题
%% chapter: 电磁感应
%% point: 验证电磁感应实验
%% method: 无
%% score: 4
%% degree: 750
%% duplicateId: 0
%% body:
一个学生用图示的电路验证楞次定律，他在实验步骤中有重要的遗漏．他的主要实验步骤如下：
\begin{steps}
	\item
	把蓄电池、开关和线圈 A 串联成一个电路；
	\item
	把电流表和线圈 B 串联成另一个电路；
	\item
	接通电流，给线圈 A 通电，并记下线圈 A 中电流的方向．把线圈 A 插入线圈 B 中，停一会儿再取出来，当线圈 A 在插入和取出过程中，以及停止运动时，观察电流表的指针有无偏移，并记下指针偏转的方向；
	\item
	改变线圈 A 中的电流方向，按步骤 3 重做实验，观察电流表的指针有无偏转，并记下指针偏转的方向。
\end{steps}

这个学生根据上述实验记录验证楞次定律，他在实验中漏掉了什么重要实验步骤？

答：\tkanswer[8cm]{查明电流表指针的偏转方向与线圈 B 中电流方向的关系}。
\onepicture[align=right, w=8cm]{figs/21.png}

%% answer: 查明电流表指针的偏转方向与线圈 B 中电流方向的关系
\jdanswer{
查明电流表指针的偏转方向与线圈 B 中电流方向的关系。
}

%% memo:
\memoanswer{
只有先查明电流表指针的偏转方向与线圈 B 中电流方向的关系，实验中才能判断电流方向与磁通量的变化的关系。

评分标准：查明电流表指针的偏转方向与线圈 B 中电流方向的关系（4 分）。
}

\item
%% number: 22
%% paperName: 1987年全国高考第 22 题 8 分
%% typeId: 5
%% type: 辨析题
%% chapter: 电路
%% point: 闭合电路欧姆定律
%% method: 无
%% score: 8
%% degree: 550
%% duplicateId: 0
%% body:
把一个“$10\UV$，$2.0\UW$”的用电器 A（纯电阻）接到某一电动势和内阻都不变的电源上，用电器 A 实际消耗的功率是 $2.0\UW$；换上另一个“$10\UV$，$5.0\UW$”的用电器 B（纯电阻）接到这一电源上，用电器 B 实际消耗的功率有没有可能反而小于 $2.0\UW$？你如果认为不可能，试说明理由。如果认为可能，试求出用电器 B 实际消耗的功率小于 $2.0\UW$ 的条件。（设电阻不随温度改变）

%% answer: 可能。r>10\sqrt{10} Ω，E>(10+2\sqrt{10}) V
\jdanswer{
可能。$r>10\sqrt{10}\UO$，$E>(10+2\sqrt{10})\UV$。
}

%% memo:
\memoanswer{
可能。条件求如下：

用 $R_{1}$ 和 $R_{2}$ 分别表示 A 和 B 的电阻（设电阻不随温度改变），由题意可知

$R_{1}=\frac{(10\UV)^{2}}{2.0\UW}$，$R_{2}=\frac{(10\UV)^{2}}{5.0\UW}$　（1）

当 A 接到电源上时，因为它正好符合额定功率 $P_{1}$ 的要求，所以有

$P_{1}=\left(\frac{E}{R_{1}+r}\right)^{2}R_{1}=2.0\UW$　（2）

当换上 B 后，B 上实际消耗的功率为

$P_{2}=\left(\frac{E}{R_{2}+r}\right)^{2}R_{2}$　（3）

根据题意，要求

$P_{2}<2.0\UW$　（4）

由（1）（2）（3）（4）式，可解得（取合理值）：

$r>10\sqrt{10}\UO$　（5），$E>(10+2\sqrt{10})\UV$　（6）

评分标准：本题 8 分。列出（1）式给 1 分，列出（2）式给 1 分，列出（3）式给 2 分，列出（4）式给 2 分。得出（5）式和（6）式的再各给 1 分。（如果考生最后给出的是开方后的四舍五入的数值，同样给分）
}

\item
%% number: 23
%% paperName: 1987年全国高考第 23 题 8 分
%% typeId: 6
%% type: 计算题
%% chapter: 运动和力
%% point: 动量守恒定律
%% method: 无
%% score: 8
%% degree: 450
%% duplicateId: 0
%% body:
如下左图所示，在水平光滑桌面上放一质量为 $M$ 的玩具小车。在小车的平台（小车的一部分）上有一质量可忽略的弹簧，一端固定在平台上，另一端用质量为 $m$ 的小球将弹簧压缩一定距离后用细线捆住。用手将小车固定在桌面上，然后烧断细线，小球就被弹出，落在车上 A 点，$OA=s$。如果小车不固定而烧断细线，球将落在车上何处？设小车足够长，球不致落在车外。
\onepicture[align=right, w=5.5cm]{figs/23.png}

%% answer: s' = \sqrt{\frac{M+m}{M}}s
\jdanswer{
$s'=\sqrt{\frac{M+m}{M}}s$
}

%% memo:
\memoanswer{
当小车固定不动时：设平台高为 $h$，小球弹出时的速度的大小为 $v$，则由平抛运动可知

$h=\frac{1}{2}gt^{2}$，$s=vt$，

联立得 $v^{2}=\frac{gs^{2}}{2h}$　（1）

当小车不固定时：设小球弹出时相对于地面的速度的大小为 $v'$，车速的大小为 $V$，由动量守恒可知

$mv'=MV$　（2）

设弹簧的弹性势能为 $E_{p}$，由能量守恒得

$E_{p}=\frac{1}{2}mv^{2}+\frac{1}{2}MV^{2}$　（3）

又 $E_{p}=\frac{1}{2}mv^{2}$，解得 $E_{p}=\frac{mgs^{2}}{4h}$。

设小球相对于小车的速度的大小为 $v''$，则

$v''=v'+V$　（4）

设小球落在车上 $A'$ 处，$OA'=s'$，则由平抛运动可知

$s'=v''\sqrt{\frac{2h}{g}}$　（5）

由（1）（2）（3）（4）（5）式可解得

$s'=\sqrt{\frac{M+m}{M}}s$　（6）

评分标准：本题 8 分。列出（1）式给 1 分，列出（2）式给 1 分，列出（3）式给 2 分，列出（4）式给 2 分，列出（5）式给 1 分，得出（6）式给 1 分。
}

\item
%% number: 24
%% paperName: 1987年全国高考第 24 题 8 分
%% typeId: 7
%% type: 作图题
%% chapter: 光学
%% point: 全反射
%% method: 无
%% score: 8
%% degree: 450
%% duplicateId: 0
%% body:
右图中 ABCD 是一个用折射率 $n=2.4$ 的透明媒质做成的四棱柱镜（图为其横截面），$\angle A=\angle C=90^{\circ}$，$\angle B=60^{\circ}$，$AB>BC$。现有平行光线垂直入射到棱镜的 AB 面上（如图示），若每个面上的反射都不能忽略，求出射光线。要求：
\begin{enumerate}
	\item
	画出所有典型光线从入射到射出的光路图。（为了图面简洁，表示光线进行方向的箭头只在棱镜外面的光线上标出即可）
	\item
	简要说明所画光路的根据，并说明每条典型光线只可能从棱镜表面的哪部分射出。
\end{enumerate}
\onepicture[align=right, w=4cm]{\includesvg{figs/24}}

%% answer: （1）有三条典型光线①②③，光路如图所示；（2）见解析
\jdanswer{
\begin{enumerate}
	\item
	有三条典型光线①，②，③，光路如图所示。

	\item
	见解析。
\end{enumerate}
}

%% memo:
\memoanswer{
（1）有三条典型光线①，②，③，光路如图所示。

\onepicture[w=4cm]{figs/24_an1.png}
\onepicture[w=4cm]{figs/24_an2.png}
\onepicture[w=4cm]{figs/24_an3.png}

（2）因为媒质的折射率 $n=2.4$，所以媒质的临界角

$\theta=\arcsin\frac{1}{n}=\arcsin\frac{1}{2.4}<30^{\circ}$。

一条典型光线是①，垂直入射到 AB 面上 BE 之间（$CE\perp AB$），部分垂直反射；部分垂直透射。到 BC 面，因入射角 $60^{\circ}$ 大于 $\theta$，发生全反射。到 DC 面，入射角 $30^{\circ}$ 仍大于 $\theta$，又发生全反射。到 AB 面，垂直入射，部分垂直射出媒质；部分垂直反射回去，根据光路可逆性，最后由原入射处射出媒质，其反射部分又重复原路。总之，光线①只能由 AB 面上 FB（$BF=BC$）间垂直射出。

一条典型光线是②，垂直入射到 AB 面上 EF 之间，部分垂直反射；部分垂直透射。到 DC 面，入射角 $30^{\circ}$ 大于 $\theta$，发生全反射。到 BC 面，入射角 $60^{\circ}$ 大于 $\theta$，全反射。到 AB 面，垂直入射，部分垂直射出媒质；部分垂直反射回去，按光路的可逆性，由原入射处射出媒质，其反射部分又重复原路。总之，光线②只能由 AB 面上 FB 间垂直射出。

一条典型光线是③，垂直入射到 AB 面上 FA 之间，部分垂直反射；部分垂直透射。到 DC 面，入射角 $30^{\circ}$ 大于 $\theta$，全反射。到 AB 面，入射角 $60^{\circ}$ 大于 $\theta$，全反射。到 BC 面，垂直入射，部分垂直射出媒质；部分垂直反射回去，按光路的可逆性，由原入射处射出媒质，其反射部分又重复原路。总之，光线③只能由 BC 面和 AB 面上 FA 间垂直射出。

评分标准：本题 8 分。（1）占 3 分，（2）占 5 分。（1）中，画光路：①占 1 分，②占 1 分，③占 1 分。每条光线的画法只要有错误或没画，就不能给该条光线的分。（2）中，解释光路：正确说明各次全反射，给 1 分；正确利用光路的可逆性，给 1 分。三条光线的射出范围，答对一条给 1 分。本题不要求画出并讨论经 E 和 F 点入射的光线的出射线。
}

\item
%% number: 25
%% paperName: 1987年全国高考第 25 题 10 分
%% typeId: 6
%% type: 计算题
%% chapter: 曲线运动
%% point: 斜抛运动
%% method: 无
%% score: 10
%% degree: 350
%% duplicateId: 0
%% body:
用大炮轰击和炮在同一水平面上的目标。当炮筒仰角为 $\varphi_{1}$ 时，着弹点比目标偏近了一段距离 $d_{1}$；当仰角为 $\varphi_{2}$ 时，着弹点又比目标偏远了一段距离 $d_{2}$。由这些已知量，求出要想正好击中目标所需的仰角 $\varphi_{0}$。设炮弹出口速率是一定的，空气阻力不计。

（本题是附加题，成绩不计入总分）

%% answer: \varphi_{0}=\frac{1}{2}\arcsin\left(\frac{d_{1}\sin2\varphi_{1}+d_{2}\sin2\varphi_{2}}{d_{1}+d_{2}}\right)
\jdanswer{
$\varphi_{0}=\frac{1}{2}\arcsin\left(\frac{d_{1}\sin2\varphi_{1}+d_{2}\sin2\varphi_{2}}{d_{1}+d_{2}}\right)$
}

%% memo:
\memoanswer{
令 $x$、$y$ 分别表示炮弹的水平位移和竖直位移，$v_{0}$ 表示其初速度，则有

$x=v_{0}\cos\varphi\cdot t$，$y=v_{0}\sin\varphi\cdot t-\frac{1}{2}gt^{2}$，

解得 $y=\frac{\sin\varphi}{\cos\varphi}x-\frac{g}{2v_{0}^{2}\cos^{2}\varphi}x^{2}$，

当炮弹落地时，$y=0$，由上两式可得水平射程

$x=\frac{v_{0}^{2}\sin2\varphi}{g}$　（1）

当仰角为 $\varphi_{1}$、$\varphi_{2}$、$\varphi_{0}$ 时，分别有

$x_{1}=\frac{v_{0}^{2}\sin2\varphi_{1}}{g}$，$x_{2}=\frac{v_{0}^{2}\sin2\varphi_{2}}{g}$，$x_{0}=\frac{v_{0}^{2}\sin2\varphi_{0}}{g}$　（2）

已知

$x_{0}-x_{1}=d_{1}$，$x_{2}-x_{0}=d_{2}$　（3）

由（2）（3）两式得

$\frac{v_{0}^{2}\sin2\varphi_{0}}{g}-\frac{v_{0}^{2}\sin2\varphi_{1}}{g}=d_{1}$，$\frac{v_{0}^{2}\sin2\varphi_{2}}{g}-\frac{v_{0}^{2}\sin2\varphi_{0}}{g}=d_{2}$，

由以上两式得 $\frac{\sin2\varphi_{0}-\sin2\varphi_{1}}{\sin2\varphi_{2}-\sin2\varphi_{0}}=\frac{d_{1}}{d_{2}}$，

解得 $\sin2\varphi_{0}=\frac{d_{1}\sin2\varphi_{1}+d_{2}\sin2\varphi_{2}}{d_{1}+d_{2}}$，

即 $\varphi_{0}=\frac{1}{2}\arcsin\left(\frac{d_{1}\sin2\varphi_{1}+d_{2}\sin2\varphi_{2}}{d_{1}+d_{2}}\right)$　（4）

评分标准：本题 10 分。列出（1）式给 4 分，列出（2）式给 2 分，列出（3）式给 2 分，得出（4）式再给 2 分。
}

\end{enumerate}

\end{document}
